\documentclass[10pt,a4paper]{amsart}
\usepackage[margin=19mm]{geometry}
\usepackage{amsmath,amssymb,mathtools}
\usepackage[scr=rsfs,cal=euler]{mathalfa}
\usepackage{xfrac}
\usepackage[unicode,hidelinks,bookmarks=false]{hyperref}
\newcommand{\ie}{i.e.\ }
\newcommand{\cf}{cf.\ }
\newcommand{\resp}{respectively\ }
\newcommand{\Subsection}{\subsection}
\providecommand{\nobreakdash}{\nobreak\hbox{-}}
\setlength{\emergencystretch}{2em}
\multlinegap0pt
\usepackage{bm}
\usepackage{bbm}
\usepackage{dsfont}
\usepackage{xspace}
\usepackage{cancel}
\usepackage{mathtools}
\usepackage{cleveref}
\usepackage{amsthm}
\makeatletter
\@ifundefined{theorem}{\newtheorem{theorem}{Theorem}}{}
\@ifundefined{lemma}{\newtheorem{lemma}[theorem]{Lemma}}{}
\@ifundefined{proposition}{\newtheorem{proposition}[theorem]{Proposition}}{}
\makeatother
\makeatletter
\newcommand{\calA}{\mathcal{A}}
\expandafter\ifx\csname customcor\endcsname\relax
\newenvironment{customcor}[1]{\renewcommand\theinnercustomcor{#1}\innercustomcor}
  {\endinnercustomcor}
\fi
\expandafter\ifx\csname customprop\endcsname\relax
\newenvironment{customprop}[1]{\renewcommand\theinnercustomprop{#1}\innercustomprop}
  {\endinnercustomprop}
\fi
\expandafter\ifx\csname customthm\endcsname\relax
\newenvironment{customthm}[1] {\renewcommand\theinnercustomthm{#1}\innercustomthm}
  {\endinnercustomthm}
\fi
\newcommand{\hroot}{\tilde{\alpha}}
\makeatother
\title[The support of Kostant's weight multiplicity formula is an o]{The support of Kostant's weight multiplicity formula is an order ideal in the weak Bruhat order}
\author{Portia X. Anderson, Esther Banaian, Melanie J. Ferreri, Owen C. Goff, Kimberly P. Hadaway, Pamela E. Harris, Kimberly J. Harry, Nicholas Mayers, Shiyun Wang,, Alexander N. Wilson}
\date{Source: arXiv:2412.16820v2 (CC BY 4.0)}
\begin{document}
\maketitle

\noindent\emph{Source and licence.} The support of Kostant's weight multiplicity formula is an order ideal in the weak Bruhat order. Version arXiv:2412.16820v2 (CC BY 4.0), distributed under the Creative Commons Attribution 4.0 International licence, as recorded in the arXiv licence metadata for this exact version and shown on the abstract page https://arxiv.org/abs/2412.16820v2. Full rights notice: licenses/arxiv-2412.16820-CC-BY-4.0.txt.\par\smallskip

\noindent\textit{Editorial scope.} Counted unit: the Lemma whose source label is \texttt{None} in the article (source0.tex lines 1658--1668, source label \texttt{None}), delivered together with the article's own complete original proof of it (source0.tex lines 1670--1924, one \texttt{proof} environment of 255 lines). No mathematics is written, repaired or re-derived: statement and proof are the article's own text line for line. Required context retained uncounted: prop:mu=-hroot (lines 912--922); lem:forbidden words (lines 938--943); Highest:1234 (lines 956--958). Every definition, display and label of those lines is retained unchanged. Omitted from this package and not redistributed: 1--911, 923--937, 944--955, 959--1657, 1669--1669, 1925--1932. Nothing inside a retained excerpt is omitted or altered: every retained body is delivered exactly as the article's lines are written, so no omission receipt of any kind accompanies this package. Citations: the retained excerpts carry no citation command at all, so no bibliography entry of the article is transcribed and no reference is redistributed. Reference adaptation: none was needed and none was made; every reference inside the delivered unit resolves inside it, so no unresolved reference or citation is produced. Printed slips kept literal: no printed word, symbol, hypothesis or number of the counted statement or of its proof was altered. Layout and class: the article's own class is \texttt{amsart} with packages including geometry, graphicx, enumerate, appendix, amsthm, xcolor, soul, hyperref, amssymb, mathtools, cleveref, tikz, thm-restate, todonotes; the wrapper is the lane's pinned \texttt{amsart} editorial wrapper at 10pt on A4, which restates the theorem-like environments the retained text needs and adds the article's own macro definitions that the retained text needs (\texttt{\textbackslash calA}, \texttt{\textbackslash hroot}), with extra wrapper packages (bm, bbm, dsfont, xspace, cancel, mathtools, cleveref). The delivered numbering is the unit's own. Assets: no retained span contains a figure environment, a graphics-inclusion command, a PDF-page inclusion, a table or a TikZ picture; no image file is copied into this package, the package declares no asset dependency and no image file is redistributed with it. Licence: version arXiv:2412.16820v2 of this article is distributed under the Creative Commons Attribution 4.0 International licence, as recorded in the arXiv licence metadata for this exact version and shown on the abstract page https://arxiv.org/abs/2412.16820v2. A full rights notice is delivered at licenses/arxiv-2412.16820-CC-BY-4.0.txt. Characters: every retained excerpt is pure ASCII, so the delivered engine, XeLaTeX with its default Latin Modern text font, reports no missing character.
\begin{proposition}\label{prop:mu=-hroot}
If $\sigma$ is one of the Weyl group elements
\begin{enumerate}
\item[(a)] $s_k$ with $1 \le k \le r$,
\item[(b)] $s_{k+1}s_{k}$ with $2 \le k \le r-2$,
\item[(c)] $s_{k}s_{k+1}$ with $2 \le k \le r-2$, 
\item[(d)] $s_{k}s_{k+1}s_{k}$ with $2 \le k \le r-2$, or
\item[(e)] $s_{k+2}s_{k}s_{k+1}$ with $2 \le k \le r-3$, 
\end{enumerate}
then $\sigma \in \calA(\tilde{\alpha},-\tilde{\alpha})$. 
\end{proposition} 

\begin{lemma}\label{lem:forbidden words}
The words 
\[s_2s_1,\qquad s_1s_2,\qquad s_{r-1}s_r, \qquad \text{and} \qquad s_rs_{r-1},\] and for any $2 \le i \le r-1$, the words
\[s_{i-1}s_is_{i+1}, \qquad s_{i}s_{i-1}s_{i+1}, \qquad \text{and} \qquad s_{i+1}s_is_{i-1}\]
are not contained in $\calA(\hroot,-\hroot)$.
\end{lemma}

\begin{lemma} \label{Highest:1234}
Let $1 \leq k \leq r-3$. The product of the four simple reflections $s_k$, $s_{k+1}$, $s_{k+2}$, and $s_{k+3}$ in any order is not contained in $\mathcal{A}(\hroot,-\hroot)$.    
\end{lemma}

\begin{lemma}
    Let $S$ be the set of elements listed in \Cref{prop:mu=-hroot}.
    Let $\sigma, \tau\in S$ be nonindependent elements.
    Then, the product $\sigma\tau$ falls into one of the following three cases:
    \begin{enumerate}
        \item $\sigma\tau\in S$,
        \item $\sigma\tau=\nu_1\nu_2\cdots\nu_m$ where $\{\nu_1,\nu_2,\ldots,\nu_m\}$ is a (possibly empty) pairwise independent subset of $S$ and \\
        $\ell(\nu_1)+\cdots+\ell(\nu_m)<\ell(\sigma)+\ell(\tau)$, or
        \item $\sigma\tau$ contains a forbidden subword listed in \Cref{lem:forbidden words} or~\Cref{Highest:1234}.
    \end{enumerate}
\end{lemma}

\begin{proof}
    We proceed via cases. Let $\sigma$ and $\tau$ come from one of the following (not necessarily distinct) collections:
\begin{enumerate}
        \item[(a)] $s_k$ with $1 \le k \le r$,
        \item[(b)] $s_{k+1}s_{k}$ with $2 \le k \le r-2$,
        \item[(c)] $s_{k}s_{k+1}$ with $2 \le k \le r-2$, 
        \item[(d)] $s_{k}s_{k+1}s_{k}$ with $2 \le k \le r-2$, or
        \item[(e)] $s_{k+2}s_{k}s_{k+1}$ with $2 \le k \le r-3$.
    \end{enumerate}
    For each pair of forms that $\sigma, \tau\in S$ can respectively take, we fix the indices of $\sigma$, then consider the range of indices of $\tau$ for which $\sigma$ and $\tau$ are not independent.

    To reduce the number of cases that need to be handled, consider the automorphism of the Weyl group $f:W\to W$ given by the conjugation $f(\sigma)=w_0\sigma w_0$ by the longest element $w_0$ of $W$. This map $f$ sends $s_i$ to $s_{r+1-i}$ and preserves length.  
    Note additionally that $f(S)=S$ and words of the form (a), (d), and (e) are sent to words of the same form. Words of the form (b) are taken to words of the form (c) by $f$ and vice versa. If $\sigma$ is a forbidden word as in \Cref{lem:forbidden words}, then so is $f(\sigma)$. 
    As the automorphism $f$ preserves words of the form (a), (d), and (e), while replacing a word of the form (c) with a word of the form (b), we can reduce the number of cases we consider.
\smallskip

\noindent \textbf{Products of the form (a)(a):}
    Let $\sigma = s_k$ for some $1 \leq k \leq r$. For $\tau = s_{j}$ with $1\leq j\leq r$ not to be independent with $\sigma$, we must have $j\in\{k-1,k,k+1\}$.
    \begin{itemize}
        \item Let $j=k-1$ for $2\leq k\leq r$. Then $\sigma\tau=s_ks_{k-1}$. If $k=2$ (resp., $k=r$), then $s_2s_1$ (resp., $s_rs_{r-1}$) is forbidden by \Cref{lem:forbidden words}. That is, $\sigma\tau$ falls into case (3). Otherwise, $3\leq k\leq r-1$ and $s_ks_{k-1}\in S$. That is, $\sigma\tau$ falls into case (1). 
        \item Let $j=k$ for $1\leq k\leq r$. Then $\sigma\tau=s_ks_k=1$, a product of an empty subset of $S$, with $0=\ell(1)=\ell(s_ks_{k})<\ell(s_k)+\ell(s_{k})=2$. That is, $\sigma\tau$ falls into case (2).
        \item Let $j=k+1$ for $1\leq k\leq r-1$. Then $\sigma\tau=s_ks_{k+1}$. If $k=1$ (resp., $k=r-1$), then $s_1s_2$ (resp., $s_{r-1}s_r$) is forbidden by \Cref{lem:forbidden words}. That is, $\sigma\tau$ falls into case (3). Otherwise, $2\leq k\leq r-2$ and $s_ks_{k+1}\in S$. That is, $\sigma\tau$ falls into case (1).
    \end{itemize}
\smallskip

\noindent \textbf{Products of the form (b)(b) (equivalent to (c)(c) via $f$):}
    Let $\sigma= s_{k+1}s_k$ for some $2\leq k\leq r-2$. For $\tau=s_{j+1}s_j$ with $2\leq j\leq r-2$ not to be independent with $\sigma$, we must have $j\in\{k-2,k-1,k,k+1,k+2\}$.
    \begin{itemize}
        \item Let $j=k-2$ for $4\leq k\leq r-2$. Then $\sigma\tau=(s_{k+1}s_k)(s_{k-1}s_{k-2})$ contains $s_ks_{k-1}s_{k-2}$, forbidden by \Cref{lem:forbidden words}. That is, $\sigma\tau$ falls into case (3).
        \item Let $j=k-1$ for $3\leq k\leq r-2$. Then $\sigma\tau=(s_{k+1}s_k)(s_ks_{k-1})=s_{k+1}s_{k-1}$ is a pairwise independent product with $2=\ell(s_{k+1})+\ell(s_{k-1})<\ell(\sigma)+\ell(\tau)=4$. That is, $\sigma\tau$ falls into case (2).
        \item Let $j=k$ for $2\leq k\leq r-2$. Then $\sigma\tau=(s_{k+1}s_k)(s_{k+1}s_k)=s_{k+1}s_ks_{k+1}s_k=s_{k+1}s_{k+1}s_ks_{k+1}=s_ks_{k+1}$, a basic allowable subword of length 2. That is, $\sigma\tau$ falls into case (2).
        \item Let $j=k+1$ for $2\leq k\leq r-3$. Then $\sigma\tau=(s_{k+1}s_k)(s_{k+2}s_{k+1})$ contains $s_{k+1}s_k s_{k+2}$ forbidden by \Cref{lem:forbidden words}. That is, $\sigma\tau$ falls into case (3).
        \item Let $j=k+2$ for $2\leq k\leq r-4$. Then $\sigma\tau=(s_{k+1}s_k)(s_{k+3}s_{k+2})=s_{k+3}s_{k+1}s_ks_{k+2}$ contains $s_{k+1}s_k s_{k+2}$ forbidden by \Cref{lem:forbidden words}. That is, $\sigma\tau$ falls into case (3).
    \end{itemize}
\smallskip

\noindent \textbf{Products of the form (d)(d):}
    Let $\sigma = s_ks_{k+1}s_k$ for some $2 \le k \le r-2$. For $\tau= s_js_{j+1}s_j$ with $2 \le j \le r-2$ not to be independent with $\sigma$, we must have $j\in\{k-2,k-1,k,k+1\}$.
    \begin{itemize}
        \item Let $j=k-2$ for $4\leq k\leq r-2$. Then $\sigma\tau=(s_ks_{k+1}s_k)(s_{k-2}s_{k-1}s_{k-2})$ contains $s_{k+1}s_ks_{k-2}s_{k-1}$, forbidden by \Cref{Highest:1234}. That is, $\sigma\tau$ falls into case (3).
        \item Let $j=k-1$ for $3\leq k\leq r-2$. Then $\sigma\tau=(s_ks_{k+1}s_k)(s_{k-1}s_ks_{k-1})$ contains $s_{k+1}s_ks_{k-1}$, forbidden by \Cref{lem:forbidden words}. That is, $\sigma\tau$ falls into case (3).
        \item Let $j=k$ for $2\leq k\leq r-2$. Then $\sigma\tau=(s_ks_{k+1}s_k)(s_ks_{k+1}s_k)=1$, a product of an empty subset of $S$, with $0=\ell(1)=\ell(s_ks_{k})<\ell(s_k)+\ell(s_{k})=2$. That is, $\sigma\tau$ falls into case (2).
        \item Let $j=k+1$ for $2\leq k\leq r-3$. Then  $\sigma\tau=(s_ks_{k+1}s_k)(s_{k+1}s_{k+2}s_{k+1})$ contains $s_ks_{k+1}s_{k+2}$, forbidden by \Cref{lem:forbidden words}. That is, $\sigma\tau$ falls into case (3).
    \end{itemize}
    \smallskip
    
\noindent \textbf{Products of the form (e)(e):}
    Let $\sigma = s_{k+2}s_ks_{k+1}$ for some $2 \le k \le r-3$. For $\tau=s_{j+2}s_js_{j+1}$ with $2 \le j \le r-3$ not to be independent with $\sigma$, we must have $j\in\{k-3,k-2,k-1,k,k+1,k+2,k+3\}$.
    \begin{itemize}
        \item Let $j=k-3$ for $5\leq k\leq r-3$. Then $\sigma\tau=(s_{k+2}s_ks_{k+1})(s_{k-1}s_{k-3}s_{k-2})$ contains $s_{k+2}s_ks_{k+1}s_{k-1}$, forbidden by \Cref{Highest:1234}. That is, $\sigma\tau$ falls into case (3).
        \item Let $j=k-2$ for $4\leq k\leq r-3$. Then $\sigma\tau=(s_{k+2}s_ks_{k+1})(s_{k}s_{k-2}s_{k-1})$ contains $s_{k+1}s_ks_{k-2}s_{k-1}$, forbidden by \Cref{Highest:1234}. That is, $\sigma\tau$ falls into case (3).
        \item Let $j=k-1$ for $3\leq k\leq r-3$. Then $\sigma\tau=(s_{k+2}s_ks_{k+1})(s_{k+1}s_{k-1}s_{k})=(s_{k+2})(s_ks_{k-1}s_k)$ is a pairwise independent product with $4=\ell(s_{k+2})+\ell(s_{k}s_{k-1}s_k)<\ell(\sigma)+\ell(\tau)=6$. That is, $\sigma\tau$ falls into case (2).
        \item Let $j=k$ for $2\leq k\leq r-3$. Then $\sigma\tau=(s_{k+2}s_ks_{k+1})(s_{k+2}s_ks_{k+1})$ contains $s_ks_{k+1}s_{k+2}$, forbidden by \Cref{lem:forbidden words}. That is, $\sigma\tau$ falls into case (3).
        \item Let $j=k+1$ for $2\leq k\leq r-4$. Then $\sigma\tau=(s_{k+2}s_ks_{k+1})(s_{k+3}s_{k+1}s_{k+2})$ contains $s_{k+2}s_ks_{k+1}s_{k+3}$, forbidden by \Cref{Highest:1234}. That is, $\sigma\tau$ falls into case (3).
        \item Let $j=k+2$ for $2\leq k\leq r-5$. Then $\sigma\tau=(s_{k+2}s_ks_{k+1})(s_{k+4}s_{k+2}s_{k+3})$ contains $s_{k+1}s_{k+4}s_{k+2}s_{k+3}$, forbidden by \Cref{Highest:1234}. That is, $\sigma\tau$ falls into case (3).
        \item Let $j=k+3$ for $2\leq k\leq r-6$. Then $\sigma\tau=(s_{k+2}s_ks_{k+1})(s_{k+5}s_{k+3}s_{k+4})=(s_{k+2}s_ks_{k+1})(s_{k+3}s_{k+5}s_{k+4})$ contains $s_{k+2}s_ks_{k+1}s_{k+3}$, forbidden by \Cref{Highest:1234}. That is, $\sigma\tau$ falls into case (3). 
    \end{itemize}
    
\smallskip

\noindent \textbf{Products of the form (a)(b) or (b)(a) (equivalent to (a)(c) or (c)(a) via $f$):}
    Let $\sigma = s_k$ for some $1 \leq k \leq r$. For $\tau = s_{j+1}s_{j}$ with $2\leq j\leq r-2$ not to be independent with $\sigma$, we must have $j\in\{k-2,k-1,k,k+1\}$.
    \begin{itemize}
        \item Let $j=k-2$ for $4\leq k\leq r$. Then \begin{itemize}
            \item $\sigma\tau=(s_k)(s_{k-1}s_{k-2})=s_ks_{k-1}s_{k-2}$ is forbidden by \Cref{lem:forbidden words}, and
            \item $\tau\sigma=(s_{k-1}s_{k-2})s_k=s_{k-1}s_{k-2}s_k$ is forbidden by \Cref{lem:forbidden words}.
        \end{itemize}
        {That is, both $\sigma\tau$ and $\tau\sigma$ fall into case (3).}
        \item Let $j=k-1$ for $3\leq k\leq r-1$. Then \begin{itemize}
            \item $\sigma\tau=(s_k)(s_k s_{k-1})=s_{k-1}\in S$, and
            \item $\tau\sigma=s_ks_{k-1}s_k\in S$.
        \end{itemize}
        {That is, both $\sigma\tau$ and $\tau\sigma$ fall into case (1).}
        \item Let $j=k$ for $2\leq k\leq r-2$. Then \begin{itemize}
            \item $\sigma\tau=(s_k)(s_{k+1}s_k)=s_k s_{k+1} s_k\in S$, and
            \item $\tau\sigma=(s_{k+1}s_k)s_k=s_{k+1}\in S$.
        \end{itemize} 
        {That is, both $\sigma\tau$ and $\tau\sigma$ fall into case (1).}
        \item Let $j=k+1$ for $1\leq k\leq r-3$. Then\begin{itemize}
            \item $\sigma\tau=(s_k)(s_{k+2}s_{k+1})=s_k s_{k+2}s_{k+1}=s_{k+2}s_k s_{k+1}\in S$ if $1<k\leq r-3$ and contains the forbidden $s_1s_2$ when $k=1$, and
            \item $\tau\sigma=(s_{k+2}s_{k+1})s_k=s_{k+2}s_{k+1}s_k$ is forbidden by \Cref{lem:forbidden words}.
        \end{itemize}
        That is, both $\sigma\tau$ and $\tau\sigma$ fall into case (3) when $k=1$, while $\sigma\tau$ falls into case (1) and $\tau\sigma$ falls into case (3) when $1<k\le r-3$.
    \end{itemize}

    
\smallskip

\noindent \textbf{Products of the form (a)(d) or (d)(a):}
    Let $\sigma=s_k$ for some $1\leq k\leq r$. For $\tau=s_js_{j+1}s_j$ with $2\leq j\leq r-2$ not to be independent with $\sigma$, we must have $j\in\{k-2, k-1, k, k+1\}$.
    \begin{itemize}
        \item Let $j=k-2$ for $4\leq k\leq r$. Then \begin{itemize}
            \item $\sigma\tau=s_k(s_{k-2}s_{k-1}s_{k-2})$ contains $s_ks_{k-1}s_{k-2}$, forbidden by \Cref{lem:forbidden words}, and
            \item $\tau\sigma=(s_{k-2}s_{k-1}s_{k-2})s_k=(s_{k-1}s_{k-2}s_{k-1})s_k$
            contains $s_{k-2}s_{k-1}s_k$, forbidden by \Cref{lem:forbidden words}.
        \end{itemize}
        That is, both $\sigma\tau$ and $\tau\sigma$ fall into case (3).
        \item Let $j=k-1$ for $3\leq k\leq r-1$. Then \begin{itemize}
            \item $\sigma\tau=s_k(s_{k-1}s_ks_{k-1})=s_k(s_ks_{k-1}s_k)=s_{k-1}s_k\in S$, and
            \item $\tau\sigma=(s_{k-1}s_ks_{k-1})s_k=(s_ks_{k-1}s_k)s_k=s_ks_{k-1}\in S$.
        \end{itemize}
        That is, both $\sigma\tau$ and $\tau\sigma$ fall into case (1).
        \item Let $j=k$ for $2\leq k\leq r-2$. Then \begin{itemize}
            \item $\sigma\tau=s_k(s_ks_{k+1}s_k)=s_{k+1}s_k\in S$, and
            \item $\tau\sigma=(s_ks_{k+1}s_k)s_k=s_ks_{k+1}\in S$.
        \end{itemize}
        That is, both $\sigma\tau$ and $\tau\sigma$ fall into case (1).
        \item Let $j=k+1$ for $1\leq k\leq r-3$. Then \begin{itemize}
            \item $\sigma\tau=s_k(s_{k+1}s_{k+2}s_{k+1})$ contains $s_ks_{k+1}s_{k+2}$, forbidden by \Cref{lem:forbidden words}, and
            \item $\tau\sigma=(s_{k+1}s_{k+2}s_{k+1})s_k$ contains $s_{k+2}s_{k+1}s_k$, forbidden by \Cref{lem:forbidden words}.
        \end{itemize}
        That is, both $\sigma\tau$ and $\tau\sigma$ fall into case (3).
    \end{itemize}
\smallskip
\noindent \textbf{Products of the form (a)(e) or (e)(a):}
    Let $\sigma=s_k$ for some $1\leq k\leq r$. For $\tau=s_{j+2}s_{j}s_{j+1}$ with $2\leq j\leq r-3$ not to be independent with $\sigma$, we must have $j\in\{k-3, k-2, k-1, k, k+1\}$.
    \begin{itemize}
        \item Let $j=k-3$ for $5\leq k\leq r$. Then \begin{itemize}
            \item $\sigma\tau= (s_{k})(s_{k-1}s_{k-3}s_{k-2})$ is forbidden by \Cref{Highest:1234}, and
            \item $\tau\sigma=(s_{k-1}s_{k-3}s_{k-2})(s_{k})$ is forbidden by \Cref{Highest:1234}.
        \end{itemize}
        That is, both $\sigma\tau$ and $\tau\sigma$ fall into case (3).
        \item Let $j=k-2$ for $4\leq k\leq r-1$. Then \begin{itemize}
            \item $\sigma\tau= (s_{k})(s_{k}s_{k-2}s_{k-1}) = s_{k-2}s_{k-1}\in S$, and
            \item $\tau\sigma= (s_{k}s_{k-2}s_{k-1})(s_{k})$ contains $s_{k-2}s_{k-1}s_k$, forbidden by \Cref{lem:forbidden words}.
        \end{itemize}
        That is, $\sigma\tau$ falls into case (1) and $\tau\sigma$ falls into case (3).
        \item Let $j=k-1$ for $3\leq k\leq r-2$. Then \begin{itemize}
            \item $\sigma\tau= (s_{k})(s_{k+1}s_{k-1}s_{k})$ contains $s_{k}s_{k-1}s_{k+1}$, forbidden by \Cref{lem:forbidden words}, and
            \item $\tau\sigma= (s_{k+1}s_{k-1}s_{k})(s_{k}) = s_{k+1}s_{k-1}$ is a product of independent elements of $S$, with $2=\ell(s_{k+1})+\ell(s_{k-1})<\ell(\tau)+\ell(\sigma)=4$.
        \end{itemize}
        That is, $\sigma\tau$ falls into case (3) and $\tau\sigma$ falls into case (2).
        \item Let $j=k$ for $2\leq k\leq r-3$. Then \begin{itemize}
            \item $\sigma\tau= (s_{k})(s_{k+2}s_{k}s_{k+1}) = s_{k+2}s_{k+1}\in S$, and
            \item $\tau\sigma= (s_{k+2}s_{k}s_{k+1})(s_{k})=s_{k+2}s_{k+1}s_ks_{k+1}$ contains $s_{k+2}s_{k+1}s_k$, forbidden by \Cref{lem:forbidden words}.
        \end{itemize}
        That is, $\sigma\tau$ falls into case (1) and $\tau\sigma$ falls into case (3).
        \item Let $j=k+1$ for $1\leq k\leq r-4$. Then \begin{itemize}
            \item $\sigma\tau= (s_{k})(s_{k+3}s_{k+1}s_{k+2})=s_{k+3}s_ks_{k+1}s_{k+2}$ contains $s_ks_{k+1}s_{k+2}$, forbidden by \Cref{lem:forbidden words}, and
            \item $\tau\sigma= (s_{k+3}s_{k+1}s_{k+2})(s_{k})$ contains $s_{k+1}s_ks_{k+2}$, forbidden by \Cref{lem:forbidden words}.
        \end{itemize}
        That is, both $\sigma\tau$ and $\tau\sigma$ fall into case (3).
    \end{itemize}
\smallskip

\noindent \textbf{Products of the form (b)(c) (equivalent to (c)(b) via $f$):}
    Let $\sigma=s_{k+1}s_k$ for some $2\leq k\leq r-2$. For $\tau=s_js_{j+1}$ with $2\leq j\leq r-2$ not to be independent with $\sigma$, we must have $j\in\{k-2,k-1,k,k+1,k+2\}$.
    \begin{itemize}
        \item Let $j=k-2$ for $4\leq k\leq r-2$. Then $\sigma\tau= (s_{k+1}s_k)(s_{k-2}s_{k-1})$ is forbidden by \Cref{Highest:1234}. That is, $\sigma\tau$ falls into case (3).
        \item Let $j=k-1$ for $3\leq k\leq r-2$. Then $\sigma\tau= (s_{k+1}s_k)(s_{k-1}s_{k})$ contains $(s_{k+1}s_ks_{k-1})$, forbidden by \Cref{lem:forbidden words}. That is, $\sigma\tau$ falls into case (3).
        \item Let $j=k$ for $2\leq k\leq r-2$. Then $\sigma\tau= (s_{k+1}s_k)(s_{k}s_{k+1})=1$, a product of an empty subset of $S$, {with $0=\ell(1)<\ell(s_{k+1}s_k)+\ell(s_{k}s_{k+1})=4$.} That is, $\sigma\tau$ falls into case (2).
        \item Let $j=k+1$ for $2\leq k\leq r-3$. Then $\sigma\tau= (s_{k+1}s_k)(s_{k+1}s_{k+2}) $ contains $s_ks_{k+1}s_{k+2}$, forbidden by \Cref{lem:forbidden words}. That is, $\sigma\tau$ falls into case (3).
        \item Let $j=k+2$ for $2\leq k\leq r-4$. Then $\sigma\tau= (s_{k+1}s_k)(s_{k+2}s_{k+3})$ contains $s_{k+1}s_ks_{k+2}$, forbidden by \Cref{lem:forbidden words}. That is, $\sigma\tau$ falls into case (3).
    \end{itemize}
    
\smallskip

\noindent \textbf{Products of the form (b)(d) or (d)(b) (equivalent to (c)(d) or (d)(c) via $f$):}
    Let $\sigma=s_{k+1}s_k$ for some $2\leq k\leq r-2$. For $\tau=s_js_{j+1}s_j$ with $2\leq j\leq r-2$ not to be independent with $\sigma$, we must have $j\in\{k-2,k-1,k,k+1,k+2\}$.
    \begin{itemize}
        \item Let $j=k-2$ for $4\leq k\leq r-2$. Then \begin{itemize}
            \item $\sigma\tau=(s_{k+1}s_k)(s_{k-2}s_{k-1}s_{k-2})$ contains $s_{k+1}s_ks_{k-2}s_{k-1}$, forbidden by \Cref{Highest:1234}, and
            \item $\tau\sigma=(s_{k-2}s_{k-1}s_{k-2})(s_{k+1}s_k)$ contains $s_{k-1}s_{k-2}s_{k+1}s_k$, forbidden by \Cref{Highest:1234}.
        \end{itemize}
        That is, both $\sigma\tau$ and $\tau\sigma$ fall into case (3).
        \item Let $j=k-1$ for $3\leq k\leq r-2$. Then \begin{itemize}
            \item $\sigma\tau=(s_{k+1}s_k)(s_{k-1}s_ks_{k-1})$ contains $s_{k+1}s_ks_{k-1}$, forbidden by \Cref{lem:forbidden words} and
            \item $\tau\sigma=(s_{k-1}s_ks_{k-1})(s_{k+1}s_k)$ contains $s_ks_{k-1}s_{k+1}$, forbidden by \Cref{lem:forbidden words}.
        \end{itemize}
        That is, both $\sigma\tau$ and $\tau\sigma$ fall into case (3).
        \item Let $j=k$ for $2\leq k\leq r-2$. Then \begin{itemize}
            \item $\sigma\tau=(s_{k+1}s_k)(s_{k}s_{k+1}s_k)=s_k\in S$, and
            \item $\tau\sigma=(s_{k}s_{k+1}s_k)(s_{k+1}s_k)=s_{k+1}\in S$.
        \end{itemize}
        That is, both $\sigma\tau$ and $\tau\sigma$ fall into case (1).
        \item Let $j=k+1$ for $2\leq k\leq r-3$. Then \begin{itemize}
            \item $\sigma\tau=(s_{k+1}s_k)(s_{k+1}s_{k+2}s_{k+1}) = s_ks_{k+1}s_ks_{k+2}s_{k+1}$ contains $s_{k+1}s_ks_{k+2}$, forbidden by \newline \Cref{lem:forbidden words}, and
            \item $\tau\sigma=(s_{k+1}s_{k+2}s_{k+1})(s_{k+1}s_k) = s_{k+1}s_{k+2}s_k$ is forbidden by \Cref{lem:forbidden words}.
        \end{itemize}
        That is, both $\sigma\tau$ and $\tau\sigma$ fall into case (3).
        \item Let $j=k+2$ for $2\leq k\leq r-4$. Then \begin{itemize}
            \item $\sigma\tau=(s_{k+1}s_k)(s_{k+2}s_{k+3}s_{k+2})$ contains $s_{k+1}s_ks_{k+2}s_{k+3}$, forbidden by \Cref{Highest:1234}, and
            \item $\tau\sigma=(s_{k+2}s_{k+3}s_{k+2})(s_{k+1}s_k)$ contains $s_{k+3}s_{k+2}s_{k+1}s_k$, forbidden by \Cref{Highest:1234}.
        \end{itemize}
        That is, both $\sigma\tau$ and $\tau\sigma$ fall into case (3).
    \end{itemize}

\smallskip
\noindent \textbf{Products of the form (b)(e) or (e)(b) (equivalent to (c)(e) or (e)(c) via $f$):}
    Let $\sigma=s_{k+1}s_k$ for some $2\leq k\leq r-2$. For $\tau=s_{j+2}s_js_{j+1}$ with $2\leq j\leq r-3$ not to be independent with $\sigma$, we must have $j\in\{k-3,k-2,k-1,k,k+1,k+2\}$.
        \begin{itemize}
        \item Let $j=k-3$ for $5\leq k\leq r-2$. Then \begin{itemize}
            \item $\sigma\tau=(s_{k+1}s_k)(s_{k-1}s_{k-3}s_{k-2})$ contains $s_ks_{k-1}s_{k-3}s_{k-2}$, forbidden by \Cref{Highest:1234}, and
            \item $\tau\sigma= (s_{k-1}s_{k-3}s_{k-2})(s_{k+1}s_k)= s_{k-3}s_{k-1}s_{k-2}s_{k+1}s_k$ contains $s_{k-1}s_{k-2}s_{k+1}s_k$, forbidden by \Cref{Highest:1234}.
        \end{itemize}
        That is, both $\sigma\tau$ and $\tau\sigma$ fall into case (3).
        \item Let $j=k-2$ for $4\leq k\leq r-2$. Then \begin{itemize}
            \item $\sigma\tau=(s_{k+1}s_k)(s_{k}s_{k-2}s_{k-1})= (s_{k+1})(s_{k-2}s_{k-1})$ is a pairwise independent product with $3=\ell(s_{k+1})+\ell(s_{k-2}s_{k-1})<\ell(\sigma)+\ell(\tau)=5$, and
            \item $\tau\sigma= (s_{k}s_{k-2}s_{k-1})(s_{k+1}s_k)$ contains $s_{k}s_{k-2}s_{k-1}s_{k+1}$, forbidden by \Cref{Highest:1234}.
        \end{itemize}
        That is, $\sigma\tau$ falls into case (2) and $\tau\sigma$ falls into case (3).
        \item Let $j=k-1$ for $3\leq k\leq r-2$. Then \begin{itemize}
            \item $\sigma\tau=(s_{k+1}s_k)(s_{k+1}s_{k-1}s_{k}) = s_ks_{k+1}s_ks_{k-1}s_k$ contains $s_{k+1}s_ks_{k-1}$, forbidden by \Cref{lem:forbidden words}, and
            \item $\tau\sigma= (s_{k+1}s_{k-1}s_{k})(s_{k+1}s_k)$ contains $s_{k-1}s_{k}s_{k+1}$, forbidden by \Cref{lem:forbidden words}.
        \end{itemize}
        That is, both $\sigma\tau$ and $\tau\sigma$ fall into case (3).
        \item Let $j=k$ for $2\leq k\leq r-3$. Then \begin{itemize}
            \item $\sigma\tau=(s_{k+1}s_k)(s_{k+2}s_{k}s_{k+1})$ contains $s_{k+1}s_ks_{k+2}$, forbidden by \Cref{lem:forbidden words}.
            \item $\tau\sigma= (s_{k+2}s_{k}s_{k+1})(s_{k+1}s_k)=s_{k+2}\in S$.
        \end{itemize}
        That is, $\sigma\tau$ falls into case (3) and $\tau\sigma$ falls into case (1).
        \item Let $j=k+1$ for $2\leq k\leq r-4$. Then \begin{itemize}
            \item $\sigma\tau=(s_{k+1}s_k)(s_{k+3}s_{k+1}s_{k+2})$ contains $s_ks_{k+3}s_{k+1}s_{k+2}$, forbidden by \Cref{Highest:1234}, and
            \item $\tau\sigma= (s_{k+3}s_{k+1}s_{k+2})(s_{k+1}s_k)$ contains $s_{k+2}s_{k+1}s_k$, forbidden by \Cref{lem:forbidden words}.
        \end{itemize}
        That is, both $\sigma\tau$ and $\tau\sigma$ fall into case (3).
    \end{itemize}
\smallskip

\noindent \textbf{Products of the form (d)(e) or (e)(d):}
    Let $\sigma=s_ks_{k+1}s_k$ for some $2\leq k\leq r-2$. For $\tau=s_{j+2}s_js_{j+1}$ with $2\leq j\leq r-3$ not to be independent with $\sigma$, we must have $j\in\{k-3,k-2,k-1,k,k+1,k+2\}$.
        \begin{itemize}
        \item Let $j=k-3$ for $5\leq k\leq r-2$. Then \begin{itemize}
            \item $\sigma\tau=(s_ks_{k+1}s_k)(s_{k-1}s_{k-3}s_{k-2})$ contains $s_ks_{k-1}s_{k-3}s_{k-2}$, forbidden by \Cref{Highest:1234}, and
            \item $\tau\sigma=(s_{k-1}s_{k-3}s_{k-2})(s_ks_{k+1}s_k)$ contains $s_{k-1}s_{k-3}s_{k-2}s_k$, forbidden by \Cref{Highest:1234}.
        \end{itemize}
        That is, both $\sigma\tau$ and $\tau\sigma$ fall into case (3).
        \item Let $j=k-2$ for $4\leq k\leq r-2$. Then \begin{itemize}
            \item $\sigma\tau=(s_ks_{k+1}s_k)(s_{k}s_{k-2}s_{k-1})=s_ks_{k+1}s_{k-2}s_{k-1}$ is forbidden by \Cref{Highest:1234}, and
            \item $\tau\sigma=(s_{k}s_{k-2}s_{k-1})(s_ks_{k+1}s_k)$ contains $s_{k-2}s_{k-1}s_ks_{k+1}$, forbidden by \Cref{Highest:1234}.
        \end{itemize}
        That is, both $\sigma\tau$ and $\tau\sigma$ fall into case (3).
        \item Let $j=k-1$ for $3\leq k\leq r-2$. Then \begin{itemize}
            \item $\sigma\tau=(s_ks_{k+1}s_k)(s_{k+1}s_{k-1}s_{k}) = (s_{k+1}s_ks_{k+1})(s_{k+1}s_{k-1}s_k) = s_{k+1}s_ks_{k-1}s_k$ contains \newline $s_{k+1}s_ks_{k-1}$, forbidden by \Cref{lem:forbidden words}, and
            \item $\tau\sigma=(s_{k+1}s_{k-1}s_{k})(s_ks_{k+1}s_k) = s_{k+1}s_{k-1}s_{k+1}s_k = s_{k-1}s_k\in S$.
        \end{itemize}
        That is, $\sigma\tau$ falls into case (3) and $\tau\sigma$ falls into case (1).
        \item Let $j=k$ for $2\leq k\leq r-3$. Then \begin{itemize}
            \item $\sigma\tau=(s_ks_{k+1}s_k)(s_{k+2}s_{k}s_{k+1})$ contains $s_{k+1}s_ks_{k+2}$, forbidden by \Cref{lem:forbidden words}, and
            \item $\tau\sigma=(s_{k+2}s_{k}s_{k+1})(s_ks_{k+1}s_k)=(s_{k+2}s_ks_{k+1})(s_{k+1}s_ks_{k+1}) = s_{k+2}s_{k+1}\in S$.
        \end{itemize}
        That is, $\sigma\tau$ falls into case (3) and $\tau\sigma$ falls into case (1).
        \item Let $j=k+1$ for $2\leq k\leq r-4$. Then \begin{itemize}
            \item $\sigma\tau=(s_ks_{k+1}s_k)(s_{k+3}s_{k+1}s_{k+2})$ contains $s_ks_{k+3}s_{k+1}s_{k+2}$, forbidden by \Cref{Highest:1234}, and
            \item $\tau\sigma=(s_{k+3}s_{k+1}s_{k+2})(s_ks_{k+1}s_k)$ contains $s_{k+3}s_{k+1}s_{k+2}s_k$, forbidden by \Cref{Highest:1234}.
        \end{itemize}
        That is, both $\sigma\tau$ and $\tau\sigma$ fall into case (3).
        \item Let $j=k+2$ for $2\leq k\leq r-5$. Then \begin{itemize}
            \item $\sigma\tau=(s_{k+1}s_{k}s_{k+1})(s_{k+4}s_{k+2}s_{k+3})$ contains $s_{k+1}s_{k+4}s_{k+2}s_{k+3}$, forbidden by \Cref{Highest:1234}, and
            \item $\tau\sigma=(s_{k+4}s_{k+2}s_{k+3})(s_{k+1}s_ks_{k+1})$ contains $s_{k+4}s_{k+2}s_{k+3}s_{k+1}$, forbidden by \Cref{Highest:1234}.
        \end{itemize}
        That is, both $\sigma\tau$ and $\tau\sigma$ fall into case (3).
    \end{itemize}
   This completes the proof.
\end{proof}

\end{document}
